% Vietnamese translation for OI Public Library
% Source-File: IOI中国国家候选队论文/2005/周源/周源.doc
\documentclass[11pt]{article}
\usepackage{oipl-vn}

\newcommand{\Oh}{\mathcal{O}}
\newcommand{\lei}{\ensuremath{\mathrm{lei}}}
\newcommand{\euro}{\ensuremath{\mathrm{euro}}}

\title{Bỏ phần thừa, giữ tinh hoa: bàn về ``phương pháp nén'' trong thi tin học}
\author{Zhou Yuan}
\date{2005}

\begin{document}
\maketitle

\origtitle{Bỏ phần thừa, giữ tinh hoa}
\sourcepath{IOI China candidate team papers / 2005 / Zhou Yuan.doc}

\renewcommand{\abstractname}{Tóm tắt}
\begin{abstract}
Trong thi tin học, ta thường gặp những bài toán có quy mô dữ liệu lớn hoặc
quan hệ giữa dữ liệu rất phức tạp; nói cách khác, lượng thông tin trong dữ
liệu đầu vào quá lớn. Tác giả gọi một cách xử lý những bài toán như vậy là
``phương pháp nén'': loại bỏ thông tin dư thừa trong dữ liệu, chỉ giữ lại phần
tinh hoa có ích cho lời giải. Phương pháp này có phạm vi ứng dụng rộng, nhưng
ở từng bài toán lại có biểu hiện rất khác nhau. Bài viết chọn một số ví dụ
tiêu biểu, rút ra điểm chung của chúng và nêu hai yếu tố then chốt để áp dụng
phương pháp nén. Phần kết luận nhấn mạnh rằng phương pháp nén kế thừa tư tưởng
quy giản, đồng thời bổ sung yếu tố ``thông tin hóa'': dùng thông tin một cách
hợp lý để làm đơn giản bài toán.
\end{abstract}

\noindent\textbf{Từ khóa:} thi tin học; phương pháp nén; thông tin dư thừa;
thông tin tinh hoa; tính nén được; quy tắc thay thế; tư tưởng quy giản; thông
tin hóa.

\section{Mở đầu}

Có lẽ nhiều bạn thấy tên gọi ``phương pháp nén'' còn lạ. Điều đó không sai,
vì đây là một thuật ngữ do tác giả đặt ra. Bài viết này sẽ dẫn dắt bạn đọc
làm quen với ý nghĩa của nó.

Khi thấy chữ ``nén'', nhiều người sẽ nghĩ tới những phần mềm quen thuộc như
WinZip hay WinRAR: chúng cố gắng giảm dung lượng tệp để phục vụ người dùng.
Đó cũng là nghĩa gốc của từ nén: dùng áp lực để giảm thể tích, kích thước,
thời gian kéo dài, mật độ hoặc nồng độ. Trong bài viết này, ``phương pháp
nén'' được hiểu là: bỏ thông tin dư thừa, giữ phần tinh hoa, từ đó giảm quy
mô bài toán. Đây chính là một cách làm rất phù hợp với thi tin học.

\section{Định nghĩa của phương pháp nén}

Khi xử lý bài toán, ta thường gặp một tập hợp $S$. Nếu quan hệ giữa các phần
tử bên trong $S$ không ảnh hưởng tới kết quả bài toán, đồng thời cũng không bị
các yếu tố bên ngoài can thiệp, thì có thể xem $S$ như một ``gói'' tách biệt
với bên ngoài. Khi đó ta bỏ qua thông tin dư thừa bên trong gói, giữ lại các
liên hệ của gói với thế giới bên ngoài, rồi nén cả gói thành một phần tử mới
để tiếp tục phân tích.

Đó là quy trình cơ bản của phương pháp nén. Lợi ích của nó là bỏ đi những chi
tiết rối rắm nhưng vô dụng bên trong gói, làm đơn giản bài toán, rồi dùng một
phương pháp tốt hơn để giải phần còn lại.

\section{Những ví dụ đơn giản}

Thực ra chúng ta không hề xa lạ với phương pháp nén. Hai ví dụ dưới đây đều
là những kỹ thuật quen thuộc.

\subsection{Bài toán đường đi ngắn nhất nhiều nguồn}

\paragraph{Mô tả bài toán.}
Trong một đồ thị có hướng trọng số dương $G=(V,E)$, cho một tập nguồn $S$.
Hãy tìm độ dài đường đi ngắn nhất từ $S$ tới mọi đỉnh còn lại. Khác với bài
toán đường đi ngắn nhất thông thường, nguồn ở đây là tất cả các đỉnh trong
$S$; khoảng cách từ $S$ tới một đỉnh $p$ bằng giá trị nhỏ nhất trong các khoảng
cách từ từng đỉnh của $S$ tới $p$.

\paragraph{Phân tích.}
Mấu chốt của bài toán là có nhiều nguồn, nhưng để đạt hiệu quả tốt ta chỉ muốn
chạy Dijkstra một lần.

Điều này rất đơn giản và có nhiều cách thực hiện. Nhìn theo phương pháp nén,
vì một đường đi ngắn nhất không cần xuất phát từ một nguồn rồi lại đi qua một
nguồn khác, quan hệ cạnh giữa các đỉnh trong $S$ không ảnh hưởng tới đáp án.
Ta có thể xem $S$ là một gói tách biệt, bỏ các thông tin dư thừa bên trong và
nén nó thành một đỉnh mới $P_S$.

Mặt khác, ta phải giữ lại các cạnh từ các đỉnh của $S$ ra bên ngoài làm cạnh
của $P_S$. Sau khi nén, ta thu được một đồ thị mới và một bài toán đường đi
ngắn nhất đơn nguồn; đó chính là việc Dijkstra cần làm.

\subsection{Bài toán đội bóng}

\paragraph{Mô tả bài toán.}
Cho đồ thị liên lạc có hướng $G=(V,E)$ của một đội bóng rổ. Nếu có cạnh
$(u,v)$ thì cầu thủ $u$ có thể kịp thời báo tin cho cầu thủ $v$. Huấn luyện
viên muốn báo một tin khẩn cấp cho toàn đội. Dựa vào đồ thị liên lạc này, ông
ấy phải trực tiếp thông báo cho ít nhất bao nhiêu cầu thủ?

\paragraph{Phân tích.}
Đồ thị liên lạc ban đầu có thể rất rối. Ta dùng phương pháp nén để bỏ bớt các
quan hệ không cần thiết.

Xét một thành phần liên thông mạnh $S$. Nếu một cầu thủ trong $S$ biết tin,
thì mọi cầu thủ khác trong $S$ cũng sẽ kịp thời biết tin. Vì vậy trạng thái
``đã biết tin'' của mọi cầu thủ trong cùng $S$ là tương đương; nếu cần, huấn
luyện viên nhiều nhất chỉ phải thông báo cho một người trong $S$. Các cạnh
bên trong $S$ vì thế không còn ý nghĩa đối với đáp án. Ta bỏ chúng, nén $S$
thành một đỉnh $P_S$, và giữ các cạnh liên lạc từ $S$ ra ngoài làm cạnh của
$P_S$.

Sau khi nén mọi thành phần liên thông mạnh, ta có một đồ thị mới không còn
chu trình. Khi đó bài toán trở nên rõ ràng: đỉnh không có bậc vào không thể
được cầu thủ khác báo tin, nên huấn luyện viên bắt buộc phải trực tiếp thông
báo cho nó. Ngược lại, trong đồ thị không chu trình, mọi đỉnh có bậc vào đều
có thể được một đỉnh không có bậc vào báo tin trực tiếp hoặc gián tiếp, nên
không cần thông báo riêng. Do đó số người tối thiểu cần thông báo bằng số
đỉnh có bậc vào bằng $0$ trong đồ thị đã nén.

Hai ví dụ trên cho thấy phương pháp nén hoạt động rất hiệu quả. Ta thường gọi
riêng dạng này trong lý thuyết đồ thị là ``co đỉnh'', nhưng ứng dụng của
phương pháp nén không chỉ giới hạn ở đồ thị. Miễn là bài toán thỏa một số
điều kiện nhất định, ta đều có thể thử tìm cách nén.

\section{Hai yếu tố then chốt}

\subsection{Tính nén được}

Tính nén được là cơ sở lý thuyết bảo đảm rằng phương pháp nén khả thi.
Muốn áp dụng phương pháp nén cho một bài toán, trước hết bài toán phải tồn
tại một hoặc nhiều loại tính nén được. Nói cách khác, trong những điều kiện
nhất định, một số tập phần tử của bài toán có thể bị nén lại. Ở hai ví dụ
trên, đó lần lượt là tập nguồn và các thành phần liên thông mạnh.

Nhìn chung, nếu liên hệ giữa các phần tử bên trong một tập $S$ không tương
tác với các phần tử bên ngoài theo cách ảnh hưởng tới kết quả bài toán, thì ta
có thể bỏ qua các liên hệ nội bộ ấy và nén $S$ thành một cá thể mới.

Trong ví dụ nhiều nguồn, quan hệ có thể đi được giữa hai nguồn bất kỳ không
ảnh hưởng tới đáp án. Vì vậy ta bỏ các liên hệ giữa các đỉnh nguồn và xét cả
tập $S$ như một nguồn duy nhất. Trong ví dụ đội bóng, mọi cầu thủ trong cùng
một thành phần liên thông mạnh luôn có trạng thái biết tin giống nhau; các
liên hệ nội bộ của họ không làm thay đổi sự thật đó, nên có thể nén cả thành
phần thành một đỉnh.

\subsection{Quy tắc thay thế}

Cần chú ý rằng nén không phải là xóa bỏ hoàn toàn. ``Nén'' gồm hai phần: phần
``ép'' cho phép bỏ tính chất cá thể và quan hệ nội bộ của tập bị nén $S$ vì
chúng là thông tin dư thừa; còn phần ``rút gọn'' phải giữ lại tinh hoa, tức
các liên hệ của $S$ với thông tin bên ngoài.

Nếu tính nén được trả lời câu hỏi vì sao có thể nén, thì quy tắc thay thế trả
lời câu hỏi nén xong dùng hình thức nào để biểu diễn phần tinh hoa còn lại.

Trong hai ví dụ đầu, quy tắc thay thế rất đơn giản: dùng một đỉnh mới thay cho
một tập nhiều đỉnh, và dùng cạnh có hướng cùng kiểu để thay cho các cạnh từ
từng đỉnh cũ ra ngoài. Dạng phần tử trước và sau khi nén hoàn toàn giống
nhau. Ở những bài phức tạp hơn, ta phải suy nghĩ kỹ hơn: sau khi nén đã giữ
lại thông tin gì, và giữ dưới dạng nào?

\subsection{Ví dụ: hệ phương trình đồng dư}

Khi giải hệ phương trình đồng dư một ẩn, ta cũng có thể xem việc gộp từng cặp
phương trình là một phép nén. Tính khả thi của phép nén là hiển nhiên; vấn đề
đáng quan tâm là quy tắc thay thế.

Giả sử cần nén hai phương trình
\[
  x \equiv a_1 \pmod{m_1}, \qquad
  x \equiv a_2 \pmod{m_2}.
\]
Ta muốn thay chúng bằng một phương trình duy nhất có cùng tập nghiệm.

Viết hai điều kiện trên dưới dạng phương trình Diophantine:
\[
  x = a_1 + m_1 t_1,\qquad
  x = a_2 + m_2 t_2.
\]
Điều này tương đương với
\[
  m_1 t_1 - m_2 t_2 = a_2-a_1.
\]
Nếu phương trình có nghiệm, ta tìm một nghiệm riêng $x_0$ của $x$, rồi tập
nghiệm chung có dạng
\[
  x = x_0 + \operatorname{lcm}(m_1,m_2)\,t,\qquad t\in\mathbb{Z}.
\]
Do đó hai phương trình ban đầu được thay bằng
\[
  x \equiv x_0 \pmod{\operatorname{lcm}(m_1,m_2)}.
\]
Đó chính là quy tắc thay thế của ví dụ này.

Tới đây ta đã nắm được hai điểm cơ bản của phương pháp nén. Trong nhiều bài
thi, chỉ cần tìm ra biểu hiện cụ thể của tính nén được và quy tắc thay thế,
ta đã tìm được điểm đặt lực để áp dụng phương pháp nén.

\section{Ứng dụng: đổi euro}

\subsection{Mô tả bài toán}

Bài toán này đến từ BOI 2003. Mỗi ngày bạn nhận được một lượng euro, có thể
dương hoặc âm. Ngân hàng cho phép bạn vào một số ngày đổi toàn bộ số euro đang
cầm sang \lei. Tỷ giá ngày thứ $i$ là
\[
  1\euro : i\lei,
\]
và mỗi lần đổi bạn phải trả thêm phí cố định $T\lei$. Tới ngày $N$ bạn bắt
buộc phải đổi hết số euro còn lại. Hãy chọn phương án để số \lei nhận được
sau ngày $N$ là lớn nhất.

Giới hạn dữ liệu:
\[
  1\le N\le 34567,\qquad 0\le T\le 34567.
\]

\subsection{Quy hoạch động ban đầu}

Gọi dãy thu nhập là $a_1,a_2,\ldots,a_N$, trong đó mỗi $a_i$ có thể âm. Một
phương án đổi tiền tương ứng với một cách chọn các điểm đổi
\[
  1\le c_1<c_2<\cdots<c_k=N.
\]
Thu nhập trong đoạn $(c_{r-1},c_r]$ được đổi ở ngày $c_r$, với $c_0=0$. Hàm
mục tiêu là
\[
  Z=\sum_{r=1}^{k}\left(c_r\sum_{j=c_{r-1}+1}^{c_r}a_j-T\right),
\]
và ta cần cực đại hóa $Z$.

Đặt $f(n)$ là lợi nhuận lớn nhất của đoạn ngày $1$ tới $n$, với $n$ là điểm
đổi cuối cùng. Khi đó
\[
  f(0)=0
\]
và
\[
  f(n)=\max_{0\le i<n}\left\{
    f(i)+n\sum_{j=i+1}^{n}a_j-T
  \right\}.
\]
Nếu đặt $S_i=\sum_{j=1}^{i}a_j$, công thức trở thành
\[
  f(n)=\max_{0\le i<n}\{f(i)+n(S_n-S_i)-T\}.
\]
Thuật toán trực tiếp có độ phức tạp $\Oh(N^2)$, không đáp ứng giới hạn dữ
liệu. Với những công thức tương tự, có thể dùng tối ưu hóa bằng bao lồi: mỗi
trạng thái quyết định $i$ được biểu diễn bởi điểm
\[
  P_i=(S_i,f(i)).
\]
Khi $S_i$ được thêm theo thứ tự đơn điệu, ta duy trì một đường gấp khúc lồi và
tìm quyết định tối ưu trong thời gian khấu hao $\Oh(1)$, thu được thuật toán
tuyến tính.

Vấn đề là ở đây $a_i$ có thể âm, nên $S_i$ không nhất thiết tăng, thậm chí có
thể dao động tùy ý. Khi tính $f(n)$, điểm mới $(S_n,f(n))$ có thể phải chèn
vào giữa đường gấp khúc hiện có, khiến việc duy trì bao lồi trở nên khó khăn.
Mâu thuẫn đã rõ: tối ưu hóa bao lồi truyền thống cần tính đơn điệu của $S_i$,
trong khi điều kiện đề bài cho phép $S_i$ thay đổi tùy ý.

\subsection{Nén để giải mâu thuẫn}

Ta quay lại ý nghĩa của đề bài. Tỷ giá ngày thứ $i$ là $1\euro:i\lei$, nên
euro càng về sau càng có giá. Nếu số euro đang cầm là dương, ta không vội đổi
vì đổi muộn hơn sẽ được nhiều hơn; nếu số euro đang cầm là âm, ta muốn bán ra
càng sớm càng tốt để chịu tổn thất ít hơn, nhưng cũng phải cân nhắc giảm số
lần đổi vì mỗi lần đổi tốn phí $T$.

Nhận xét cảm tính này dẫn tới định lý sau.

\paragraph{Định lý 4.1.}
Giả sử dãy $\{a_i\}$ có một đoạn liên tiếp
\[
  a_i,a_{i+1},\ldots,a_j,\qquad i<j,
\]
sao cho các tổng tiền tố
\[
  a_i,\quad a_i+a_{i+1},\quad \ldots,\quad
  a_i+a_{i+1}+\cdots+a_{j-1}
\]
đều dương, còn tổng cả đoạn
\[
  a_i+a_{i+1}+\cdots+a_j
\]
không dương. Khi đó tồn tại một phương án tối ưu trong đó sau khi nhận
$a_i$ ở ngày $i$, ta không đổi ngay, mà tiếp tục giữ qua các ngày
$i+1,\ldots,j$ rồi mới xét có đổi hay không.

\paragraph{Chứng minh.}
Dùng phản chứng và xét tình huống đơn giản nhất: giả sử một phương án tối ưu
có đúng một điểm đổi $k$ trong đoạn $[i,j)$.

Gọi $x$ là lượng euro còn lại từ điểm đổi trước tới hết ngày $i-1$. Nếu $x$
dương, thì đến ngày $k$ lượng euro còn lại càng dương hơn. Bỏ điểm đổi $k$ và
gộp nó vào điểm đổi sau không chỉ làm euro được đổi ở tỷ giá cao hơn, mà còn
tiết kiệm một khoản phí $T$.

Nếu $x$ âm, ta dời điểm đổi $k$ về ngày $i-1$. Phần euro âm được bán sớm hơn,
nên tổn thất tính theo \lei nhỏ hơn; còn đoạn $[i,k]$ có tổng tiền tố dương
được gộp vào điểm đổi sau, cũng không làm kết quả xấu đi. Trong cả hai trường
hợp, ta đều cải thiện hoặc không làm giảm lợi nhuận, mâu thuẫn với giả thiết
về tính tối ưu của phương án ban đầu. Trường hợp có nhiều điểm đổi trong đoạn
được xử lý tương tự.

Định lý trên cho thấy bài toán có tính nén được. Nếu từ một vị trí bắt đầu,
các tổng lũy tích vẫn dương, ta chắc chắn chưa cần đổi; chỉ tới khi tổng này
trở nên không dương mới cần xét tiếp. Những ngày trong đoạn đó là tương đương
về mặt quyết định đổi tiền: trong một phương án tối ưu, chúng được xử lý ở
cùng một điểm đổi. Vì vậy ta có thể nén chúng thành một điểm.

Cụ thể, nếu đoạn đầu tiên bị nén là $a_1,\ldots,a_p$, đặt
\[
  b_1=a_1+a_2+\cdots+a_p,\qquad t_1=p.
\]
Ta dùng cặp $(b_1,t_1)$ thay cho $p$ khoản thu nhập ban đầu. Giá trị $b_1$
cho biết lượng euro sau nén, còn $t_1$ cho biết thời điểm sớm nhất có thể đổi
khoản tiền nén này. Tiếp tục xử lý từ $a_{p+1}$, ta nhận được
\[
  (b_1,t_1),(b_2,t_2),\ldots,(b_M,t_M).
\]
Nếu tới ngày cuối cùng mà lượng euro còn lại vẫn dương, ta giữ nó tới ngày
$N$, vì đó là thời điểm có tỷ giá tốt nhất.

Trên dãy đã nén, ta vẫn viết một quy hoạch động tương tự. Đặt
\[
  S'_m=\sum_{r=1}^{m}b_r.
\]
Vì mỗi $b_r\le 0$ sau khi nén theo quy tắc trên, dãy $S'_m$ không tăng. Điều
kiện đơn điệu cần cho tối ưu hóa bằng bao lồi đã xuất hiện. Do đó ta áp dụng
phương pháp trong phụ lục để tính quy hoạch động trên dãy nén trong thời gian
$\Oh(M)$, cộng với thời gian tiền xử lý $\Oh(N)$. Tổng độ phức tạp là
\[
  \Oh(N+M)=\Oh(N).
\]

Ví dụ này cho thấy sức mạnh của phương pháp nén: điều kiện ``$a_i$ có thể âm
dương tùy ý'' phá vỡ tính đơn điệu cần cho tối ưu hóa quy hoạch động, nhưng
việc nén các đoạn có tổng tiền tố dương đã biến dãy mới thành một dãy có tổng
tiền tố đơn điệu, mở đường cho thuật toán tuyến tính.

\section{Ứng dụng: hợp nhất dãy}

\subsection{Mô tả bài toán}

Bài toán này được cải biên từ ZOJ p1794, \textit{Merging Sequence Problem}.
Cho $K$ dãy, độ dài lần lượt là $n_1,n_2,\ldots,n_K$. Mỗi số trong các dãy có
thể âm. Ta cần trộn $K$ dãy này thành một dãy mới $S$ như sau: mỗi bước chọn
một dãy đầu vào không rỗng, lấy phần tử đầu tiên của dãy đó đặt vào cuối $S$,
rồi xóa phần tử ấy khỏi dãy đầu vào. Lặp lại cho tới khi tất cả dãy đều rỗng.

Gọi độ dài của $S$ là
\[
  N=n_1+n_2+\cdots+n_K,
\]
và các tổng tiền tố của $S$ là
\[
  \mathrm{Sum}_0,\mathrm{Sum}_1,\ldots,\mathrm{Sum}_N.
\]
Hãy tìm một cách trộn để giá trị lớn nhất trong các tổng tiền tố này nhỏ nhất
có thể. Giới hạn là
\[
  K\le N\le 10^5.
\]

\subsection{Quan sát điểm có thể nén}

Thoạt nhìn, bài toán dường như chỉ có một quy hoạch động rất lớn, chẳng hạn
kiểu $\Oh(KN^K)$, hoàn toàn không thực tế. Ta thử phân tích có ý thức theo
phương pháp nén.

Giả sử trong một dãy đầu vào có một đoạn con liên tiếp
\[
  u_1,u_2,\ldots,u_p.
\]
Nếu chứng minh được rằng khi quá trình trộn đã chọn $u_1$ đưa vào $S$, thì ta
có thể tiếp tục chọn liên tiếp tới $u_p$ mà không làm mất khả năng đạt nghiệm
tối ưu, tức là đoạn $u$ xuất hiện liên tiếp trong một nghiệm tối ưu nào đó,
thì đoạn $u$ có thể nén. Khi ấy quan hệ vị trí nội bộ của các phần tử trong
$u$ không còn cần xét riêng nữa.

Quy tắc thay thế cần giữ lại hai thông tin. Thứ nhất, một tổng tiền tố bên
trong đoạn $u$ có thể trở thành tổng tiền tố lớn nhất của toàn bộ $S$; chỉ
đỉnh tương đối lớn nhất trong đoạn là đáng quan tâm. Thứ hai, tổng của cả đoạn
$u$ sẽ ảnh hưởng tới các tổng tiền tố phía sau. Vì vậy ta dùng một cặp
\[
  (a,b)
\]
để mô tả đặc trưng của một đoạn đã nén: $a\ge 0$ là đỉnh tương đối của đoạn,
$b\le 0$ là phần hiệu chỉnh sao cho $a+b$ bằng tổng của cả đoạn. Nói cách
khác, đoạn được thay bằng một ``đỉnh đơn'' gồm bước đi lên $a$ rồi đi xuống
$b$.

Một cặp $(a,b)$ như vậy, với $b\le 0$, rõ ràng có thể coi là một đơn vị nén.
Ta gọi nó là một \term{cặp}. Nếu $a+b>0$ thì cặp là dương, nếu $a+b<0$ thì
cặp là âm, và nếu $a+b=0$ thì gọi là cặp không.

\subsection{Giai đoạn nén thứ nhất}

\paragraph{Định lý 5.1.}
Giả sử trong một dãy đầu vào tồn tại đoạn liên tiếp
\[
  u_1,u_2,\ldots,u_p
\]
có các tổng tiền tố
\[
  Q_1,Q_2,\ldots,Q_p,
\]
trong đó $Q_1,\ldots,Q_{p-1}$ đều không âm, còn $Q_p<0$. Khi có thể chọn
$u_1$ đưa vào $S$, thì hoặc tạm thời chưa chọn, hoặc phải chọn liên tiếp từ
$u_1$ tới $u_p$. Vì vậy đoạn này có thể nén.

Để chứng minh định lý, trước hết xét bổ đề sau.

\paragraph{Bổ đề 5.1.1.}
Nếu trong một dãy đầu vào có hai cặp liên tiếp
\[
  (a_1,b_1),(a_2,b_2),
\]
trong đó cặp đầu không âm và cặp sau âm, thì có thể nén hai cặp này thành một
cặp mới.

\paragraph{Chứng minh.}
Giả sử trong $S$ có đoạn gồm ba phần:
\[
  (a_1,b_1),\quad (x,y),\quad (a_2,b_2),
\]
trong đó $(x,y)$ là một đoạn nằm giữa hai cặp cần xét.

Nếu $x+y\le 0$, ta đổi thứ tự thành
\[
  (x,y),\quad (a_1,b_1),\quad (a_2,b_2).
\]
Vì $(a_1,b_1)$ không âm, đỉnh tuyệt đối của $x$ không tăng; vì $x+y\le 0$,
đỉnh tuyệt đối của $a_1$ cũng không tăng.

Nếu $x+y>0$, ta đổi thứ tự thành
\[
  (a_1,b_1),\quad (a_2,b_2),\quad (x,y).
\]
Vì $(a_2,b_2)$ là cặp âm, đỉnh của $(x,y)$ giảm; đồng thời đỉnh của
$(a_2,b_2)$ cũng không tăng. Trong cả hai trường hợp, ta đều có thể đưa hai
cặp đang xét về vị trí liên tiếp mà không làm tăng đỉnh lớn nhất của $S$.
Do đó hai cặp có thể nén.

Từ bổ đề này suy ra định lý 5.1. Ban đầu xem mỗi số riêng lẻ như một cặp đặc
biệt. Miễn là $p$ số chưa bị nén hết thành một cặp, cặp đầu tiên vẫn là dương;
vì tổng cả đoạn là âm, luôn tồn tại một cặp âm ở phía sau. Ta chọn một cặp
dương đứng trước một cặp âm và dùng bổ đề để nén chúng. Lặp lại cho tới khi
chỉ còn một cặp, ta chứng minh được cả đoạn có thể nén.

Vì vậy, với mỗi dãy đầu vào, ta quét từ trái sang phải, cộng dồn tổng. Nếu
tổng hiện tại vẫn dương thì tiếp tục; khi tổng trở thành âm, ta nén đoạn vừa
xét thành một cặp, đặt tổng về $0$ rồi xử lý tiếp.

Sau giai đoạn này, phần đầu của mỗi dãy đầu vào được biến thành nhiều cặp âm
$(a_i,b_i)$: mỗi số dương đều đi kèm một số âm có trị tuyệt đối lớn hơn.
Phần đuôi còn lại là một đoạn mà mọi tổng tiền tố đều dương. Có thể chứng minh
rằng trong một nghiệm tối ưu, toàn bộ các đoạn đuôi dương này sẽ nằm ở nửa sau
của dãy $S$. Nếu nhìn ngược các đoạn ấy, bài toán của chúng tương đương với
phần đầu đã xử lý. Vì vậy ta chỉ cần tập trung vào bài toán gồm các cặp âm.

\subsection{Giai đoạn nén thứ hai}

\paragraph{Định lý 5.2.}
Nếu trong một dãy đầu vào có hai cặp âm liên tiếp
\[
  (a_1,b_1),(a_2,b_2),
\]
và đỉnh nằm ở cặp thứ nhất, tức là
\[
  a_1\ge a_1+b_1+a_2,
\]
thì hai cặp này chắc chắn có thể xuất hiện liên tiếp trong một nghiệm tối ưu,
do đó có thể nén.

\paragraph{Chứng minh.}
Giả sử trong một phương án tối ưu có đoạn
\[
  (a_1,b_1),\quad (x,y),\quad (a_2,b_2).
\]
Đỉnh lớn nhất của đoạn này là
\[
  \max\{a_1,\ a_1+b_1+x,\ a_1+b_1+x+y+a_2\}.
\]
Nếu đổi thành
\[
  (a_1,b_1),\quad (a_2,b_2),\quad (x,y),
\]
đỉnh lớn nhất trở thành
\[
  \max\{a_1,\ a_1+b_1+a_2,\ a_1+b_1+a_2+b_2+x\}.
\]
Do $a_2+b_2\le 0$ và $a_1\ge a_1+b_1+a_2$, giá trị sau không lớn hơn giá trị
trước. Vậy có một phương án tối ưu trong đó hai cặp đang xét đứng liên tiếp.

\paragraph{Hệ quả 5.2.1.}
Nếu trong một dãy đầu vào có nhiều cặp âm liên tiếp
\[
  (a_1,b_1),(a_2,b_2),(a_3,b_3),\ldots
\]
và đỉnh của toàn đoạn luôn nằm ở cặp đầu tiên, tức là
\[
  a_1\ge a_1+b_1+a_2,\quad
  a_1\ge a_1+b_1+a_2+b_2+a_3,\quad \ldots,
\]
thì cả đoạn cặp này có thể nén.

Dựa trên hệ quả đó, ta nén tất cả các đoạn thỏa điều kiện giai đoạn hai. Sau
khi xử lý, mỗi dãy mới có dấu của các số dương và âm xen kẽ, trị tuyệt đối
tăng dần, hay nói cách khác đỉnh tương đối của các cặp trong cùng một dãy là
đơn điệu tăng.

\subsection{Thuật toán tham lam}

Sau hai giai đoạn nén, bài toán có dạng rất gọn: có $K$ dãy, mỗi dãy gồm các
cặp có tổng âm, và đỉnh tương đối của các cặp trong từng dãy tăng dần. Cần
trộn chúng thành một dãy $S$ sao cho tổng tiền tố lớn nhất nhỏ nhất.

Vì mỗi cặp có tổng âm, đặt cặp có đỉnh tương đối cao càng muộn càng làm giảm
đỉnh tuyệt đối của nó. Đây là nguyên tắc tham lam. Do đỉnh tương đối trong
mỗi dãy đã tăng dần, ta chỉ cần trộn $K$ dãy cặp theo thứ tự tăng của đỉnh
tương đối, giống như thao tác merge nhiều dãy đã sắp xếp. Ở mỗi bước, chọn
cặp có đỉnh tương đối nhỏ nhất trong số các cặp đang đứng đầu các dãy, đưa
nó vào $S$, rồi cập nhật hàng đợi ưu tiên.

Dùng một heap, mỗi lần chọn mất $\Oh(\log K)$, nên toàn bộ thuật toán chạy
trong thời gian
\[
  \Oh(N\log K).
\]
Như vậy phương pháp nén không chỉ tìm ra một thuật toán đa thức cho bài toán,
mà còn đưa tới một thuật toán rất tốt.

Ví dụ này cũng cho thấy đôi khi một lần nén chưa đủ. Sau giai đoạn thứ nhất,
các dãy đã bớt hỗn loạn nhưng chưa đủ điều kiện để tham lam. Giai đoạn thứ hai
tiếp tục bỏ thông tin dư thừa, giữ lại đúng phần có ích, làm xuất hiện tính
đơn điệu cần cho lời giải cuối cùng.

\section{Kết luận}

Hãy nhìn lại các ví dụ chính. Ở bài đường đi ngắn nhất nhiều nguồn, ta nén tập
nguồn thành một nguồn mới, biến bài toán nhiều nguồn thành bài toán đơn
nguồn. Ở bài đội bóng, ta nén mỗi thành phần liên thông mạnh thành một đỉnh,
biến đồ thị liên lạc tùy ý thành một đồ thị có hướng không chu trình. Ở bài
hợp nhất dãy, nhờ hai loại tính nén được khác nhau, các đoạn con phù hợp đều
được thay bằng một ``đỉnh đơn'' gồm hai số $(a,b)$; từ đó $K$ dãy ban đầu sau
khi biến đổi thể hiện tính đơn điệu đẹp, cho phép dùng tham lam.

Trong các ví dụ ấy, phương pháp nén đều thay một đối tượng phức tạp bằng một
đối tượng đơn giản hơn, nhờ đó làm đơn giản bài toán. Đây chính là một tư
tưởng kinh điển trong toán học: tư tưởng quy giản. Khi giải phương trình, ta
thường dùng đặt ẩn phụ để biến một dạng khó thành một dạng quen thuộc hơn;
phương pháp nén cũng có tinh thần tương tự.

Tuy nhiên, đặt ẩn phụ thông thường chủ yếu là biến đổi hình thức và bảo toàn
tương đương. Phương pháp nén thì khác. Hai yếu tố của nó, tính nén được và
quy tắc thay thế, đều xoay quanh việc sử dụng thông tin. Tính nén được thừa
nhận rằng trong một tập có thông tin dư thừa, tức là quan hệ giữa các phần tử
không còn ý nghĩa đối với đáp án. Quy tắc thay thế đưa ra cách thao tác cụ
thể để giữ lại phần tinh hoa cần thiết.

Ta đang sống trong một xã hội bùng nổ thông tin. Thu nhận thông tin nhanh và
chính xác đã không còn là điều khó nhất; điều khó hơn là tránh bị nhiễu bởi
thông tin vô ích. Thi tin học, ở một mức độ nào đó, cũng dạy ta cách sử dụng
thông tin hợp lý và hiệu quả. Phương pháp nén kế thừa tư tưởng quy giản cổ
điển, đồng thời thêm vào yếu tố thông tin hóa: bỏ thông tin dư thừa để biến
phức tạp thành đơn giản, biến khó thành dễ.

Vì vậy, khi gặp một bài toán có quá nhiều điều kiện, lượng thông tin quá lớn
và chưa thấy đường xử lý, hãy thử nghĩ tới phương pháp nén. Mục tiêu của nó
chính là:
\[
  \text{bỏ phần thừa, giữ tinh hoa.}
\]

\appendix
\section{Ghi chú về tối ưu hóa độ dốc trong bài đổi euro}

Xét công thức quy hoạch động
\[
  f(n)=\max_{0\le i<n}\{f(i)+n(S_n-S_i)-T\}.
\]
Khi $n$ đã cố định, đặt
\[
  g_i(n)=f(i)+n(S_n-S_i)-T.
\]
Với hai quyết định $i,j$ và $S_i>S_j$, quyết định $i$ tốt hơn $j$ khi
\[
  f(i)-f(j)>n(S_i-S_j),
\]
hay tương đương
\[
  \frac{f(i)-f(j)}{S_i-S_j}>n.
\]
Nếu xem mỗi quyết định là một điểm $P_i=(S_i,f(i))$, thì việc so sánh hai
quyết định trở thành so sánh độ dốc của đoạn thẳng nối hai điểm với giá trị
$n$. Các điểm lõm không bao giờ trở thành quyết định tối ưu, nên chỉ cần duy
trì một bao lồi thích hợp. Khi thứ tự thêm điểm theo hoành độ là đơn điệu,
việc chèn điểm và lấy quyết định tối ưu có thể làm trong thời gian khấu hao
$\Oh(1)$.

\section{Phụ lục chương trình và dữ liệu}

Bản gốc có kèm đề tiếng Anh của BOI 2003 Euro và các chương trình minh họa
cho hai ví dụ chính. Theo quy ước của bản dịch này, các phụ lục chương trình
không được dịch lại. Nội dung thuật toán tương ứng đã được trình bày trong
phần phân tích.

Với chương trình cho bài hợp nhất dãy, tệp vào có tên
\texttt{MergeSequence.in}. Dòng đầu chứa số nguyên $K$, số dãy cần trộn. Mỗi
trong $K$ dòng tiếp theo mô tả một dãy: số đầu tiên là $n_i$, số lượng phần tử
của dãy, theo sau là $n_i$ số theo đúng thứ tự trong dãy đó.

\section{Tài liệu tham khảo}

\begin{itemize}
  \item Wu Wenhu, Wang Jiande, \textit{Hướng dẫn thi Olympic Tin học: thuật
  toán và lập trình đồ thị bằng Pascal}.
  \item Wu Wenhu, Wang Jiande, \textit{Phân tích và thiết kế chương trình
  thuật toán thực dụng}.
  \item Zhou Yuan, \textit{Bàn về ứng dụng tư tưởng kết hợp số và hình trong
  thi tin học}, luận văn đội tuyển tập huấn Olympic Tin học năm 2004.
  \item Zhejiang University Online Judge: \url{http://acm.zju.edu.cn}.
\end{itemize}

\end{document}
